% --- Archon physics formalization source begin ---
% archon:physics
% archon:covers PhyXMiniProblems/problem_phyx_mini_0391.lean
% archon:source-report reports/phyx_mini/problem_phyx_mini_0391.source.json

\chapter{Physics problem phyx\_mini\_0391}
\label{ch:PhyXMiniProblems_problem_phyx_mini_0391}

\paragraph{Problem source.}
\#\# Physical scenario

A piece of experimental apparatus, figure, is located where \$g = 9.5\$ m/s\$\textasciicircum{}2\$ and the temperature is \$5\textasciicircum{}\{\textbackslash{}circ\}C\$. Air flow inside the apparatus is determined by measuring the pressure drop across an orifice with a mercury manometer showing a height difference of 200 mm.

\#\# Question

What is the pressure drop in kPa?

\#\# Answer choices

- A: A: 490kPa

- B: B: 25.84kPa

- C: C: 154kPa

- D: D: 10.2kPa

\#\# Figure caption (auxiliary; use the image as primary evidence)

The image is a diagram depicting a manometer setup for measuring pressure difference. It consists of the following components:

1. **U-tube Manometer**: 
   - The manometer is U-shaped and partially filled with a blue-colored fluid, likely indicating a liquid like mercury or water.
   - The liquid is at different levels in the two arms of the U-tube.

2. **Airflow**: 
   - There is a horizontal channel at the top, indicating airflow from left to right, labeled with the word "Air."
   - The arrows on either side of the air channel represent the direction of the airflow.

3. **Probe or Tube**:
   - A vertical tube extends from the U-tube into the airflow channel, indicating the point where pressure is measured.

4. **Gravitational Force**:
   - An arrow labeled “g” is pointing downward, indicating the direction of the gravitational force affecting the liquid in the manometer.

The diagram likely illustrates how the pressure difference between the airflow and the atmospheric or another reference pressure causes the liquid to deflect, providing a measure of this pressure difference.

\#\# Dataset metadata

Temperature and Heat Transfer; Physical Model Grounding Reasoning, Multi-Formula Reasoning

\paragraph{Recorded answer/context.}
B: B: 25.84kPa

\paragraph{Figure/image path.}
/root/proposal\_for\_physic/science-mango/phyx\_mini\_run/phyx\_data/test\_image/391.png

\paragraph{Formalization target.}
create a compiling Lean file with sorry bodies at `PhyXMiniProblems/problem\_phyx\_mini\_0391.lean`. The Lean declarations must preserve the physical quantities, dimensions or dimensional roles, figure labels, governing-law hypotheses, and final relation expressed by this problem.
Use Mathlib/Physlib names found through LeanExplore where available. If a physics API is missing, introduce faithful local abstractions rather than scalar placeholder aliases.

\begin{theorem}[Physics formalization target]
\label{thm:physics:phyx_mini_0391:target}
\lean{PhyXMiniProblems.ProblemPhyXMini0391.problem_phyx_mini_0391}
\uses{def:physics:phyx-mini-0391:phyxminiproblems-problemphyxmini0391-pressureinkilopascals, def:physics:phyx-mini-0391:phyxminiproblems-problemphyxmini0391-orificemanometersetup, def:physics:phyx-mini-0391:phyxminiproblems-problemphyxmini0391-matchesprimaryfigure, def:physics:phyx-mini-0391:phyxminiproblems-problemphyxmini0391-matchesproblemdata, def:physics:phyx-mini-0391:phyxminiproblems-problemphyxmini0391-usesmercurypropertydataatfivedegrees, def:physics:phyx-mini-0391:phyxminiproblems-problemphyxmini0391-satisfiesorificemanometerlaws, def:physics:phyx-mini-0391:phyxminiproblems-problemphyxmini0391-isuniquematchingdisplayedpressure}
The assigned autoformalize agent should translate this physics problem into Lean declarations in the covered file, with theorem and lemma proof bodies written as `by sorry`.
\end{theorem}
\begin{proof}
This is an autoformalization task, not a proof task. Produce statements that can later be proved by the physics prover without weakening the physical model.
\end{proof}

% --- Archon declaration topology begin ---
\section{Declaration topology}
\begin{definition}[mass Density Dimension]
  \label{def:physics:phyx-mini-0391:phyxminiproblems-problemphyxmini0391-massdensitydimension}
  \lean{PhyXMiniProblems.ProblemPhyXMini0391.massDensityDimension}
  The physical dimension of mass density, M L⁻³
\end{definition}
\begin{proof}
  This is the declared model definition of mass Density Dimension.
\end{proof}
\begin{definition}[acceleration Dimension]
  \label{def:physics:phyx-mini-0391:phyxminiproblems-problemphyxmini0391-accelerationdimension}
  \lean{PhyXMiniProblems.ProblemPhyXMini0391.accelerationDimension}
  The physical dimension of acceleration, L T⁻²
\end{definition}
\begin{proof}
  This is the declared model definition of acceleration Dimension.
\end{proof}
\begin{definition}[Length Quantity]
  \label{def:physics:phyx-mini-0391:phyxminiproblems-problemphyxmini0391-lengthquantity}
  \lean{PhyXMiniProblems.ProblemPhyXMini0391.LengthQuantity}
  A nonnegative physical length, independent of the readout unit
\end{definition}
\begin{proof}
  This is the declared model definition of Length Quantity.
\end{proof}
\begin{definition}[Temperature Quantity]
  \label{def:physics:phyx-mini-0391:phyxminiproblems-problemphyxmini0391-temperaturequantity}
  \lean{PhyXMiniProblems.ProblemPhyXMini0391.TemperatureQuantity}
  A nonnegative absolute temperature, independent of the readout unit
\end{definition}
\begin{proof}
  This is the declared model definition of Temperature Quantity.
\end{proof}
\begin{definition}[Mass Density Quantity]
  \label{def:physics:phyx-mini-0391:phyxminiproblems-problemphyxmini0391-massdensityquantity}
  \lean{PhyXMiniProblems.ProblemPhyXMini0391.MassDensityQuantity}
  \uses{def:physics:phyx-mini-0391:phyxminiproblems-problemphyxmini0391-massdensitydimension}
  A nonnegative physical mass density
\end{definition}
\begin{proof}
  This is the declared model definition of Mass Density Quantity.
\end{proof}
\begin{definition}[Acceleration Quantity]
  \label{def:physics:phyx-mini-0391:phyxminiproblems-problemphyxmini0391-accelerationquantity}
  \lean{PhyXMiniProblems.ProblemPhyXMini0391.AccelerationQuantity}
  \uses{def:physics:phyx-mini-0391:phyxminiproblems-problemphyxmini0391-accelerationdimension}
  A nonnegative physical acceleration magnitude
\end{definition}
\begin{proof}
  This is the declared model definition of Acceleration Quantity.
\end{proof}
\begin{definition}[length Readout]
  \label{def:physics:phyx-mini-0391:phyxminiproblems-problemphyxmini0391-lengthreadout}
  \lean{PhyXMiniProblems.ProblemPhyXMini0391.lengthReadout}
  \uses{def:physics:phyx-mini-0391:phyxminiproblems-problemphyxmini0391-lengthquantity}
  Read a physical length in a selected length unit
\end{definition}
\begin{proof}
  This is the declared model definition of length Readout.
\end{proof}
\begin{definition}[density Readout]
  \label{def:physics:phyx-mini-0391:phyxminiproblems-problemphyxmini0391-densityreadout}
  \lean{PhyXMiniProblems.ProblemPhyXMini0391.densityReadout}
  \uses{def:physics:phyx-mini-0391:phyxminiproblems-problemphyxmini0391-massdensityquantity}
  Read a mass density in a selected mass unit per selected length unit cubed
\end{definition}
\begin{proof}
  This is the declared model definition of density Readout.
\end{proof}
\begin{definition}[acceleration Readout]
  \label{def:physics:phyx-mini-0391:phyxminiproblems-problemphyxmini0391-accelerationreadout}
  \lean{PhyXMiniProblems.ProblemPhyXMini0391.accelerationReadout}
  \uses{def:physics:phyx-mini-0391:phyxminiproblems-problemphyxmini0391-accelerationquantity}
  Read an acceleration in selected length and time units
\end{definition}
\begin{proof}
  This is the declared model definition of acceleration Readout.
\end{proof}
\begin{definition}[length In Meters]
  \label{def:physics:phyx-mini-0391:phyxminiproblems-problemphyxmini0391-lengthinmeters}
  \lean{PhyXMiniProblems.ProblemPhyXMini0391.lengthInMeters}
  \uses{def:physics:phyx-mini-0391:phyxminiproblems-problemphyxmini0391-lengthquantity, def:physics:phyx-mini-0391:phyxminiproblems-problemphyxmini0391-lengthreadout}
  Meter readout of a physical length
\end{definition}
\begin{proof}
  This is the declared model definition of length In Meters.
\end{proof}
\begin{definition}[length In Millimeters]
  \label{def:physics:phyx-mini-0391:phyxminiproblems-problemphyxmini0391-lengthinmillimeters}
  \lean{PhyXMiniProblems.ProblemPhyXMini0391.lengthInMillimeters}
  \uses{def:physics:phyx-mini-0391:phyxminiproblems-problemphyxmini0391-lengthquantity, def:physics:phyx-mini-0391:phyxminiproblems-problemphyxmini0391-lengthreadout}
  Millimeter readout of a physical length
\end{definition}
\begin{proof}
  This is the declared model definition of length In Millimeters.
\end{proof}
\begin{definition}[temperature In Kelvins]
  \label{def:physics:phyx-mini-0391:phyxminiproblems-problemphyxmini0391-temperatureinkelvins}
  \lean{PhyXMiniProblems.ProblemPhyXMini0391.temperatureInKelvins}
  \uses{def:physics:phyx-mini-0391:phyxminiproblems-problemphyxmini0391-temperaturequantity}
  Kelvin readout of an absolute temperature
\end{definition}
\begin{proof}
  This is the declared model definition of temperature In Kelvins.
\end{proof}
\begin{definition}[temperature In Degrees Celsius]
  \label{def:physics:phyx-mini-0391:phyxminiproblems-problemphyxmini0391-temperatureindegreescelsius}
  \lean{PhyXMiniProblems.ProblemPhyXMini0391.temperatureInDegreesCelsius}
  \uses{def:physics:phyx-mini-0391:phyxminiproblems-problemphyxmini0391-temperaturequantity, def:physics:phyx-mini-0391:phyxminiproblems-problemphyxmini0391-temperatureinkelvins}
  Celsius readout obtained from the coherent Kelvin readout
\end{definition}
\begin{proof}
  This is the declared model definition of temperature In Degrees Celsius.
\end{proof}
\begin{definition}[density In Kilograms Per Cubic Meter]
  \label{def:physics:phyx-mini-0391:phyxminiproblems-problemphyxmini0391-densityinkilogramspercubicmeter}
  \lean{PhyXMiniProblems.ProblemPhyXMini0391.densityInKilogramsPerCubicMeter}
  \uses{def:physics:phyx-mini-0391:phyxminiproblems-problemphyxmini0391-massdensityquantity, def:physics:phyx-mini-0391:phyxminiproblems-problemphyxmini0391-densityreadout}
  Kilogram-per-cubic-meter readout of a mass density
\end{definition}
\begin{proof}
  This is the declared model definition of density In Kilograms Per Cubic Meter.
\end{proof}
\begin{definition}[acceleration In Meters Per Second Squared]
  \label{def:physics:phyx-mini-0391:phyxminiproblems-problemphyxmini0391-accelerationinmeterspersecondsquared}
  \lean{PhyXMiniProblems.ProblemPhyXMini0391.accelerationInMetersPerSecondSquared}
  \uses{def:physics:phyx-mini-0391:phyxminiproblems-problemphyxmini0391-accelerationquantity, def:physics:phyx-mini-0391:phyxminiproblems-problemphyxmini0391-accelerationreadout}
  Meter-per-second-squared readout of an acceleration
\end{definition}
\begin{proof}
  This is the declared model definition of acceleration In Meters Per Second Squared.
\end{proof}
\begin{definition}[pressure In Pascals]
  \label{def:physics:phyx-mini-0391:phyxminiproblems-problemphyxmini0391-pressureinpascals}
  \lean{PhyXMiniProblems.ProblemPhyXMini0391.pressureInPascals}
  Pascal readout of a pressure
\end{definition}
\begin{proof}
  This is the declared model definition of pressure In Pascals.
\end{proof}
\begin{definition}[pressure In Kilopascals]
  \label{def:physics:phyx-mini-0391:phyxminiproblems-problemphyxmini0391-pressureinkilopascals}
  \lean{PhyXMiniProblems.ProblemPhyXMini0391.pressureInKilopascals}
  \uses{def:physics:phyx-mini-0391:phyxminiproblems-problemphyxmini0391-pressureinpascals}
  Kilopascal readout of a pressure difference
\end{definition}
\begin{proof}
  This is the declared model definition of pressure In Kilopascals.
\end{proof}
\begin{definition}[Pressure Tap]
  \label{def:physics:phyx-mini-0391:phyxminiproblems-problemphyxmini0391-pressuretap}
  \lean{PhyXMiniProblems.ProblemPhyXMini0391.PressureTap}
  The sides of the orifice at which air pressure is sampled
\end{definition}
\begin{proof}
  This is the declared model definition of Pressure Tap.
\end{proof}
\begin{definition}[Manometer Arm]
  \label{def:physics:phyx-mini-0391:phyxminiproblems-problemphyxmini0391-manometerarm}
  \lean{PhyXMiniProblems.ProblemPhyXMini0391.ManometerArm}
  The two arms visible in the U-tube drawing
\end{definition}
\begin{proof}
  This is the declared model definition of Manometer Arm.
\end{proof}
\begin{definition}[Apparatus Fluid]
  \label{def:physics:phyx-mini-0391:phyxminiproblems-problemphyxmini0391-apparatusfluid}
  \lean{PhyXMiniProblems.ProblemPhyXMini0391.ApparatusFluid}
  Fluids relevant to the apparatus
\end{definition}
\begin{proof}
  This is the declared model definition of Apparatus Fluid.
\end{proof}
\begin{definition}[Figure Text Label]
  \label{def:physics:phyx-mini-0391:phyxminiproblems-problemphyxmini0391-figuretextlabel}
  \lean{PhyXMiniProblems.ProblemPhyXMini0391.FigureTextLabel}
  Literal text labels visible in the supplied bitmap
\end{definition}
\begin{proof}
  This is the declared model definition of Figure Text Label.
\end{proof}
\begin{definition}[Airflow Direction]
  \label{def:physics:phyx-mini-0391:phyxminiproblems-problemphyxmini0391-airflowdirection}
  \lean{PhyXMiniProblems.ProblemPhyXMini0391.AirflowDirection}
  Direction of the arrows drawn in the air passage
\end{definition}
\begin{proof}
  This is the declared model definition of Airflow Direction.
\end{proof}
\begin{definition}[Orifice Manometer Figure]
  \label{def:physics:phyx-mini-0391:phyxminiproblems-problemphyxmini0391-orificemanometerfigure}
  \lean{PhyXMiniProblems.ProblemPhyXMini0391.OrificeManometerFigure}
  \uses{def:physics:phyx-mini-0391:phyxminiproblems-problemphyxmini0391-lengthquantity, def:physics:phyx-mini-0391:phyxminiproblems-problemphyxmini0391-manometerarm, def:physics:phyx-mini-0391:phyxminiproblems-problemphyxmini0391-figuretextlabel, def:physics:phyx-mini-0391:phyxminiproblems-problemphyxmini0391-airflowdirection}
  The geometric and qualitative evidence visible in the supplied bitmap. The height difference is a dimensionful quantity but receives its numerical 200 mm readout separately from the problem statement
\end{definition}
\begin{proof}
  This is the declared model definition of Orifice Manometer Figure.
\end{proof}
\begin{definition}[Orifice Manometer Setup]
  \label{def:physics:phyx-mini-0391:phyxminiproblems-problemphyxmini0391-orificemanometersetup}
  \lean{PhyXMiniProblems.ProblemPhyXMini0391.OrificeManometerSetup}
  \uses{def:physics:phyx-mini-0391:phyxminiproblems-problemphyxmini0391-lengthquantity, def:physics:phyx-mini-0391:phyxminiproblems-problemphyxmini0391-temperaturequantity, def:physics:phyx-mini-0391:phyxminiproblems-problemphyxmini0391-massdensityquantity, def:physics:phyx-mini-0391:phyxminiproblems-problemphyxmini0391-accelerationquantity, def:physics:phyx-mini-0391:phyxminiproblems-problemphyxmini0391-pressuretap, def:physics:phyx-mini-0391:phyxminiproblems-problemphyxmini0391-manometerarm, def:physics:phyx-mini-0391:phyxminiproblems-problemphyxmini0391-apparatusfluid, def:physics:phyx-mini-0391:phyxminiproblems-problemphyxmini0391-orificemanometerfigure}
  Independent physical quantities and apparatus roles. In particular, the pressure drop is an independent pressure quantity: it is not defined from the recorded answer or from the hydrostatic formula
\end{definition}
\begin{proof}
  This is the declared model definition of Orifice Manometer Setup.
\end{proof}
\begin{definition}[Matches Primary Figure]
  \label{def:physics:phyx-mini-0391:phyxminiproblems-problemphyxmini0391-matchesprimaryfigure}
  \lean{PhyXMiniProblems.ProblemPhyXMini0391.MatchesPrimaryFigure}
  \uses{def:physics:phyx-mini-0391:phyxminiproblems-problemphyxmini0391-orificemanometersetup}
  Evidence read directly from the primary bitmap phyx\_data/test\_image/391.png
\end{definition}
\begin{proof}
  This is the declared model definition of Matches Primary Figure.
\end{proof}
\begin{definition}[Matches Problem Data]
  \label{def:physics:phyx-mini-0391:phyxminiproblems-problemphyxmini0391-matchesproblemdata}
  \lean{PhyXMiniProblems.ProblemPhyXMini0391.MatchesProblemData}
  \uses{def:physics:phyx-mini-0391:phyxminiproblems-problemphyxmini0391-lengthinmillimeters, def:physics:phyx-mini-0391:phyxminiproblems-problemphyxmini0391-accelerationinmeterspersecondsquared, def:physics:phyx-mini-0391:phyxminiproblems-problemphyxmini0391-orificemanometersetup}
  Textual apparatus roles and numerical readouts supplied in the problem. The two pressure taps lie on opposite sides of the orifice; the higher upstream pressure depresses the mercury in the left arm in the pictured orientation
\end{definition}
\begin{proof}
  This is the declared model definition of Matches Problem Data.
\end{proof}
\begin{definition}[Uses Mercury Property Data At Five Degrees]
  \label{def:physics:phyx-mini-0391:phyxminiproblems-problemphyxmini0391-usesmercurypropertydataatfivedegrees}
  \lean{PhyXMiniProblems.ProblemPhyXMini0391.UsesMercuryPropertyDataAtFiveDegrees}
  \uses{def:physics:phyx-mini-0391:phyxminiproblems-problemphyxmini0391-temperatureindegreescelsius, def:physics:phyx-mini-0391:phyxminiproblems-problemphyxmini0391-densityinkilogramspercubicmeter, def:physics:phyx-mini-0391:phyxminiproblems-problemphyxmini0391-orificemanometersetup}
  Mercury property-table calibration at the stated operating temperature. The density value is a material-data readout, not the requested pressure answer
\end{definition}
\begin{proof}
  This is the declared model definition of Uses Mercury Property Data At Five Degrees.
\end{proof}
\begin{definition}[Satisfies Orifice Manometer Laws]
  \label{def:physics:phyx-mini-0391:phyxminiproblems-problemphyxmini0391-satisfiesorificemanometerlaws}
  \lean{PhyXMiniProblems.ProblemPhyXMini0391.SatisfiesOrificeManometerLaws}
  \uses{def:physics:phyx-mini-0391:phyxminiproblems-problemphyxmini0391-lengthinmeters, def:physics:phyx-mini-0391:phyxminiproblems-problemphyxmini0391-densityinkilogramspercubicmeter, def:physics:phyx-mini-0391:phyxminiproblems-problemphyxmini0391-accelerationinmeterspersecondsquared, def:physics:phyx-mini-0391:phyxminiproblems-problemphyxmini0391-pressureinpascals, def:physics:phyx-mini-0391:phyxminiproblems-problemphyxmini0391-orificemanometersetup}
  The apparatus-specific governing laws, expressed in coherent SI readouts. The second field is the U-tube hydrostatic balance Δp = ρ g h under the stated heavy-manometer-liquid approximation
\end{definition}
\begin{proof}
  This is the declared model definition of Satisfies Orifice Manometer Laws.
\end{proof}
\begin{definition}[Answer Choice]
  \label{def:physics:phyx-mini-0391:phyxminiproblems-problemphyxmini0391-answerchoice}
  \lean{PhyXMiniProblems.ProblemPhyXMini0391.AnswerChoice}
  The four answer choices in the problem statement
\end{definition}
\begin{proof}
  This is the declared model definition of Answer Choice.
\end{proof}
\begin{definition}[displayed Pressure In Kilopascals]
  \label{def:physics:phyx-mini-0391:phyxminiproblems-problemphyxmini0391-displayedpressureinkilopascals}
  \lean{PhyXMiniProblems.ProblemPhyXMini0391.displayedPressureInKilopascals}
  \uses{def:physics:phyx-mini-0391:phyxminiproblems-problemphyxmini0391-answerchoice}
  Displayed pressure-drop values, in kilopascals
\end{definition}
\begin{proof}
  This is the declared model definition of displayed Pressure In Kilopascals.
\end{proof}
\begin{definition}[recorded Dataset Answer]
  \label{def:physics:phyx-mini-0391:phyxminiproblems-problemphyxmini0391-recordeddatasetanswer}
  \lean{PhyXMiniProblems.ProblemPhyXMini0391.recordedDatasetAnswer}
  \uses{def:physics:phyx-mini-0391:phyxminiproblems-problemphyxmini0391-answerchoice}
  Dataset metadata; this declaration is not used as a theorem premise
\end{definition}
\begin{proof}
  This is the declared model definition of recorded Dataset Answer.
\end{proof}
\begin{definition}[Matches Displayed Pressure]
  \label{def:physics:phyx-mini-0391:phyxminiproblems-problemphyxmini0391-matchesdisplayedpressure}
  \lean{PhyXMiniProblems.ProblemPhyXMini0391.MatchesDisplayedPressure}
  \uses{def:physics:phyx-mini-0391:phyxminiproblems-problemphyxmini0391-pressureinkilopascals, def:physics:phyx-mini-0391:phyxminiproblems-problemphyxmini0391-orificemanometersetup, def:physics:phyx-mini-0391:phyxminiproblems-problemphyxmini0391-answerchoice, def:physics:phyx-mini-0391:phyxminiproblems-problemphyxmini0391-displayedpressureinkilopascals}
  A choice matches the pressure drop derived for the setup
\end{definition}
\begin{proof}
  This is the declared model definition of Matches Displayed Pressure.
\end{proof}
\begin{definition}[Is Unique Matching Displayed Pressure]
  \label{def:physics:phyx-mini-0391:phyxminiproblems-problemphyxmini0391-isuniquematchingdisplayedpressure}
  \lean{PhyXMiniProblems.ProblemPhyXMini0391.IsUniqueMatchingDisplayedPressure}
  \uses{def:physics:phyx-mini-0391:phyxminiproblems-problemphyxmini0391-orificemanometersetup, def:physics:phyx-mini-0391:phyxminiproblems-problemphyxmini0391-answerchoice, def:physics:phyx-mini-0391:phyxminiproblems-problemphyxmini0391-matchesdisplayedpressure}
  A choice is the unique displayed value matching the derived pressure drop
\end{definition}
\begin{proof}
  This is the declared model definition of Is Unique Matching Displayed Pressure.
\end{proof}
% --- Archon declaration topology end ---
% --- Archon physics formalization source end ---
